.mt9
.mb7
.po8
.ll65
.he
.bp38
.ig
CHAPTER 16
O P E R A T I O N A L	O V E R V I E W
.sp2
.ce
SETUP
.pp5
System setup involves four steps:
.he Operations: Chapter 16
.sp
.in10
.ti-4
1.  Make two copies of the SSG distribution disk.
.sp
.ti-4
2.  Define your CRT to the Display Management System.
.sp
.ti-4
3.  Set the system parameters.
.sp
.ti-4
4.  Build the item file.
.qi
.pp
To make the copies of the distribution disk you need two formatted
disks with the CP/M operating system SYSGEN'ed on (see the CP/M Features
and Facilities manual on how to do this).  These disks should also
hold the CP/M PIP and STAT commands, as well as the CBASIC runtime
interpreter CRUN2.  After you PIP all the Inventory Control System files
onto your disks, store the SSG distribution disk safely away; use one
of your newly created disks as a system master and one as a work disk.
If you need to make new or additional work disks, make them from the
system master disk, not the SSG distribution disk.
.pp
The Display Management System comes ready to run on a Hazeltine 1500
CRT device; if you have a Hazeltine 1500 you just erase the DEF type
files and the CRTDEF.INT program from your two disks
(not the SSG distribution disk!) and continue.	 If you have one
of the CRT's listed in the following chapter the procedure involves
renaming the file that corresponds to your CRT to CRT.DEF; you then
erase the rest of the DEF files and the CRTDEF.INT program and begin.
If you have neither a Hazeltine 1500, nor any of the CRT's listed in the
next chapter, you define your CRT to the Display Management System by
running the CRTDEF.INT program and answering its requests.
.pp
To set the system parameters you call up the Inventory Control System
by the command CRUN2 IY, set the GENERAL PARAMETERS according to your
needs, set the AUDIT PARAMETERS, set the ITEM FILE PARAMETERS, and
the SORT PARAMETERS.  At first time start-up the system takes you
immediately to the PARAMETER ENTRY program where you set the general
and audit parameters, and then to the section of the ITEM FILE RECOPY
program where you set the item file parameters.  Once these are set, the
system menu comes up and you should choose
the CREATE SORT PARAMETER FILES program.
.pp
To build the item file you call the ITEM ENTRY program through the system
menu.  If no valid records exist on the file the program goes right into
the adding mode; if records do exist on the file you must give the
add records command (CTRL-A).  Enter the item number, then the information
for the the item.  Enter ESCAPE at any time to to place the information in
the file and begin adding the next item.
.sp2
.ce
NORMAL PROCESSING
.pp
Normal processing includes adding to stock when goods arrive, depleting
stock when they leave, generating the reports, and, if you wish auditability,
printing the transaction proofs.
.pp
To add to stock you select the ADD TO STOCK program from the menu, access
the item by item number, relative record number, or sequentially by a next-
record or previous record command, and
enter the quantity added and its cost.	The program updates the quantities
on order and back-ordered by subtracting the amount added from both values.
You can then change the order dates or adjust the other
order information.  ESCAPE enters the transaction and
begins the next additions to stock.  ESCAPE once more
returns you to the system menu.
.pp
To record depletions from stock select the menu program DEPLETE STOCK.	Access
the part as for adding to stock and enter the quantity depleted.  The program
automatically updates the quantity in stock.
.pp
If you've chosen the auditability option (at the AUDIT PARAMETERS screen of
the PARAMETER ENTRY program), and want both additions and depletions
audited, you must erase the current audit file before going from
add to stock to deplete stock, or from deplete stock to add to stock.
This is because the audit file	holds additions or depletions, but not
both kinds of transactions at the same time.
To ensure accuracy and completeness you should print a transaction proof
before you erase the audit file.  If you've chosen to audit additions only,
or depletions only, there is no need to erase the file, but you should
proof and erase it regularly so it does not grow unnecessarily large.
.pp
To generate a report, select the desired reporting program from the system
menu, choose the report options available, set up your printer, and print.
.pp
When you add new items to your stock they must be added to the master item
file through the ITEM ENTRY program.  If the addition of the item means
the file is no longer ordered by item number, a sort will be required.
How you run the sort depends on how many item files you have, how many
disk drives you have, and the size of the item files.  Chapter 22 gives
full details.
.pp
If you wish to change the distribution of the items among the item files,
or when you want to physically remove deleted records from the file, you
run the program ITEM FILE RECOPY from the system menu, changing disks as
directed by the program.
.pp
If the number of item files changes, or if you want to alter your sorting
procedures, you will need to change the sort parameters through the CREATE
SORT PARAMETER FILES program.
.he
.bp
.ig
CHAPTER 17
S Y S T E M    I N S T A L L A T I O N
.pp5
The first step in installing the Inventory Control System is to make
two system disks from the SSG distribution disk.
.he Operations: Chapter 17
.pp
Place a freshly formatted disk containing the CP/M operating system,
the CBASIC runtime interpreter CRUN2, and the CP/M PIP and STAT commands
on drive A.  Place the (write-protected)
SSG distribution disk on drive B and hit the
RESET switch on your machine.  Type the command,
.sp
.ti20
A>PIP A:=B:*.*[OV]
.sp
.ad
and follow with RETURN.  When the PIP is complete, remove the distribution
disk from drive B and replace it with a formatted disk with the
CP/M operating system SYSGEN'ed onto it.  Reset the machine, then
give the command,
.sp
.ti20
A>PIP B:=A:*.*[OV]
.pp
The second step in installing the Inventory Control System is to
define the characteristics of your CRT to the Display Management
System.
.pp
If you have the Hazeltine 1500 CRT, the Inventory Control
System will run immediately, just erase all files with the filetype
(or suffix) of DEF and the CRTDEF.INT program from your work disks.
.pp
If you have any of the following CRT's you must use the CP/M REN command
to rename the DEF file	corresponding to your CRT to the filename CRT.DEF:
.cp12
.sp
.li
      CRT			      DEF FILE
   -------------------------------------------------------
   ADM 3A			      ADM3A.DEF
   BEEHIVE 150			      B150.DEF
   BEEHIVE 152			      B152.DEF
   BEEHIVE 157			      B157.DEF
   CROMEMCO 3100		      CRO3100.DEF
   DYNABYTE NAKED TERMINAL	      DYNABYTE.DEF
   SOROC IQ 120 		      SOROC120.DEF
   IMSAI VIOC			      VIOC.DEF
.bp
.ad
Type the command,
.sp
.ti20
A>REN A:CRT.DEF=A:nnnnnnnn.ttt[v]
.sp
.ad
where nnnnnnnn.ttt stands for the filename and filetype of your CRT
definition file.  Erase the remaining DEF files from your work disk
(not the SSG distribution disk!), but do not erase the CRT.DEF file.
Erase also the CRTDEF.INT program.
.pp
If you have neither the Hazeltine 1500 terminal, nor any of the CRT's
listed above, see Chapter 28 for instructions on how to define your
CRT to the DMS by running the CRTDEF.INT program.
.pp
The Inventory Control System will always expect to find the CRT.DEF
file on the currently logged drive.
.pp
You can now call up the Inventory Control System and begin setting
the system parameters and building the item file, the last steps remaining
in the system setup sequence.
.he
.bp
.ig
CHAPTER 18
I N T E R A C T I N G	W I T H   T H E    S Y S T E M
.pp5
At first time startup
the Inventory Control System creates a default parameter file,
goes directly into the PARAMETER ENTRY program, then into the ITEM FILE
RECOPY program (to set the item file parameters), and then to the system
menu.  Once the system is installed, the command
.he Operations: Chapter 18
.sp
.ti20
A>CRUN2 IY
.sp
.ad
brings up the system menu immediately.	All functions except for
sorting under certain circumstances (see Chapter 22), can be performed
by choosing programs from the system menu (see Chapter 27 for a reproduction
of the system menu).
.pp
Entry programs use formatted screen displays with prompts for information
and bracketed fields to receive the data or
parameter values you enter.  RETURN moves the CRT cursor forward
along a predetermined route
through the fields on the screen.  CTRL-B backs the cursor up one field in the
DMS screen.
.pp
If your entry does not fill the field completely you must hit RETURN to
go on to the next field.  Whatever you type into the field becomes the
current value for that field, regardless of what other characters may
still remain in the field.  For example, if the current value in the
Company Name field is [Ajax Enterprises 	], and your new entry
is "Ajax Inc.", the field would look like this before you hit the
RETURN key: [Ajax Inc.rprises	      ].  After hitting RETURN the field
would show [Ajax Inc.----------------].
.pp
To end any Inventory Control System function press the ESCAPE key.
ESCAPE takes you one "layer" out of the program.  For example,
while adding items to the file you hit ESCAPE to begin adding the next
item.  ESCAPE again takes you out of the adding mode but leaves
you at the ITEM ENTRY screen where you can examine, change, or delete
other file records.  ESCAPE once more removes you from the ITEM ENTRY
screen and program altogether and returns the system menu.  ESCAPE at the
system menu ends the Inventory Control System.	The procedure for leaving
the ADD TO STOCK and DEPLETE STOCK programs is similar, though you need not
hit ESCAPE between depletion transactions.
Typing the word
STOP works similarly at non-DMS controlled requests.
.pp
To cancel a transaction such as entering a new item into the file, adding
to stock, or depleting stock, issue the command CTRL-X at any field on the
DMS screen.  Should a DMS screen become garbled for any reason, the command
CTRL-R refreshes it.
.pp
To back up one request when not in a Display Management System
screen (such as at the report option requests) you should enter an up-arrow
(^) followed by RETURN.
.pp
To back the cursor up one space in a DMS screen field, use the command
CTRL-H.  At non-DMS requests, use the DELETE command.
.he
.bp
.ig
CHAPTER 19
S E T T I N G	S Y S T E M   P A R A M E T E R S
A N D	D R I V E   A L L O C A T I O N
.pp5
The procedure for setting system parameters at initial startup differs
from the procedure for setting the parameters once the system is installed.
When no parameter files or valid inventory records exist, after asking for
the "system drive", the system goes immediately to the PARAMETER ENTRY
program where you can examine or change
the general and audit parameters.  After
these are set, the system moves on to the section of the ITEM FILE RECOPY
program that allows you to set or change the item file parameters (see
Chapter 23). After the item file parameters are set, the system menu
comes up.  You should choose the CREATE SORT PARAMETER FILES program, and
then the ITEM ENTRY program.  If no valid records exist on the item file
the program goes automatically into the adding mode, with the cursor in the
item number field. You then enter the item number
and other information for that item.
.he Operations: Chapter 19
.pp
After first time startup, to change the parameters
you call the appropriate program from the system menu directly.
All fields for parameter (and data) entry are edited to exclude double
quotation marks (") and invalid
control characters, and to ensure the correctness
of the entry (e.g., alphabetic characters not allowed in the vendor number
field).
.sp
.ce
GENERAL PARAMETERS
.pp
Set the general parameters through the menu function PARAMETER ENTRY.  This
brings up the PARAMETER ENTRY MENU.  Choose GENERAL PARAMETERS and make
the desired entries or changes subject to the restrictions listed below:
.sp
.in15
.ti-10
COMPANY NAME:  Maximum length is 60 characters
.sp
.ti-10
CURRENT PERIOD:  Accepts week, month, or quarter only.
.sp
.ti-10
TO DATE PERIOD:  Accepts month, quarter, or year only.
.sp
.ti-10
DATE FORMAT:  Accepts MM/DD/YY or DD/MM/YY.
.sp
.ti-10
PAGE WIDTH:  Accepts any value from 1 to 999.  Since the Inventory
Control System requires a 132 column printer, the PAGE WIDTH parameter
control the placement of the report headings.  A value less than 132
moves the heading to the left, a value greater than 132 moves it to
the right.
.sp
.ti-10
LINES PER PAGE:  Accepts any value from 3 to 999.  The default is 60
lines, or 10 inches, measured from the top of the first printed line.
.sp
.ti-10
AVERAGE VALUATION:  When set to NO, the value of the inventory must be
calculated manually, outside the system, and the INVENTORY VALUE field
updated through the ITEM ENTRY program if you wish a valuation report.
.sp
.ti-10
ARRAY SIZE:  This value depends on the amount of Random Access Memory
(RAM) available.  The larger the array, the faster the record access,
although unless you expect to access many items during a run of the system
a larger array size does not substantially speed up record access.
Use the chart below to determine the right size for your computer (memory
sizes are CP/M configured):
.sp
.cp6
.li
		   RAM		      ARRAY SIZE
		------------------------------------
		   45K			  100
		   48K			  175
		   56K			  500
		   64K			 1000
.sp
.ti-10
CONSOLE BELL:  Sounds the console "bell" when errors occur.
.sp
.ti-10
DEBUGGING:  Leave set to NO.
.qi
.sp3
.ce
AUDIT PARAMETERS
.pp
Set the audit parameters through the PARAMETER ENTRY program.  When the
PARAMETER ENTRY MENU comes up, choose AUDIT PARAMETERS.  Enter the
disk drive where the audit file is to reside (A-D), and then set the
ADDITIONS AUDIT USED and DEPLETIONS AUDIT USED values as desired.
For no auditability, set both to NO.
.bp
.ce
ITEM FILE PARAMETERS
.pp
Set or change the item file parameters through the menu program ITEM FILE
RECOPY.  When asked if you wish to change the parameters answer YES.  A NO
answer tells the program you don't wish to change the parameters, but want
to remove deleted records from the file.  If you anticipate only one
item file, enter the drive location (A-D).  If more than one, enter the drive
location for each, and the lowest item number on item file #2 and item file #3.
Item numbers may contain only the letters A through Z, the numbers 0 through
9, dashes, or slashes.
.sp3
.ce
SORT PARAMETERS
.pp
For each item file you order, there is a corresponding sort parameter file.
If you change the number of item files, or if they grow to a size that
makes the current set of sort parameters unworkable, you will need to change
the values in the sort parameter files.  After calling the menu program
CREATE SORT FILE PARAMETERS, set the workfile and output drive locations,
and indicate whether you want to be able to switch disks to receive the sorted
output.  The program assumes the input drive is the same drive as that
specified for the item file.  The sort parameter files are named (if three
item files exist) PART1.SRT, PART2.SRT, and PART3.SRT.
.sp3
.ce
DISK DRIVE ALLOCATION
.sp1
.pp
The Inventory Control System always
expects to find the program files (those with
the filetype of INT) and QSORT
on the currently logged drive.	The two parameter
files IYPR1.101 AND IYPR2.101 are placed on the drive specified as the
"system" drive at system startup.  The system drive must be either drive A or
B only.  If they are found to be on a different
drive than the one specified as the system drive, the program assumes that
that drive is now the system drive.  The sort parameter files PART1.SRT,
PART2.SRT, and PART3.SRT are also expected to be on the system drive,
unless they are placed on a special sorting disk (see Chapter 22). The
audit file drive is established through the PARAMETER ENTRY program and can
be on any drive location.  The CRT.DEF file (if used) is expected to reside on
the currently logged drive.
.pp
The chart below suggests the best file allocation for one to three
item files on two to four disk drives:
.bp
.li
  2 DRIVE SYSTEM
  --------------------------------------------------------------
  1 Item File:	       A: INT FILES
			  QSORT
			  SYSTEM DRIVE: IYPR1
					IYPR2
					SRT FILES
			  AUDIT FILE (IYAUD)
		       B: ITEM FILE #1 (IYPRT)

  2 Item Files:        A: INT FILES
			  QSORT
			  ITEM FILE #1 (IYPRT)
		       B: SYSTEM DRIVE: IYPR1
					IYPR2
					SRT FILES
			  AUDIT FILE (IYAUD)
			  ITEM FILE #2 (IYPRT)
.br
.sp6
.li
  3 DRIVE SYSTEM
  --------------------------------------------------------------
  1 Item File:	       A: INT FILES
			  QSORT
		       B: SYSTEM DRIVE: IYPR1
					IYPR2
					SRT FILES
			  AUDIT FILE (IYAUD)
		       C: ITEM FILE #1 (IYPRT)

  2 Item Files:        A: INT FILES
			  QSORT
			  ITEM FILE #1 (IYPRT)
		       B: SYSTEM DRIVE: IYPR1
					IYPR2
					SRT FILES
			  AUDIT FILE (IYAUD)
		       B: ITEM FILE #1 (IYPRT)
		       C: ITEM FILE #2 (IYPRT)

.bp
.li
  3 Item Files:        A: INT FILES
			  QSORT
			  ITEM FILE #1 (IYPRT)
		       B: SYSTEM DRIVE: IYPR1
					IYPR2
					SRT FILES
			  AUDIT FILE (IYAUD)
			  ITEM FILE #2 (IYPRT)
		       C: ITEM FILE #3 (IYPRT)

.br
.sp6
.li
  4 DRIVE SYSTEM
  -------------------------------------------------------------
  1 Item File:	       A: INT FILES
			  QSORT
		       B: SYSTEM DRIVE: IYPR1
					IYPR2
					SRT FILES
			  AUDIT FILE (IYAUD)
		       C: ITEM FILE #1 (IYPRT)
		       D: not used, or AUDIT FILE

  2 Item Files:        A: INT FILES
			  QSORT
		       B: SYSTEM DRIVE: IYPR1
					IYPR2
					SRT FILES
			  AUDIT FILE (IYAUD)
		       C: ITEM FILE #1 (IYPRT)
		       D: ITEM FILE #2 (IYPRT)

  3 Item Files:        A: INT FILES
			  QSORT
			  SYSTEM DRIVE: IYPR1
					IYPR2
					SRT FILES
			  AUDIT FILE (IYAUD)
		       B: ITEM FILE #1 (IYPRT)
		       C: ITEM FILE #2 (IYPRT)
		       D: ITEM FILE #3 (IYPRT)

.he
.bp
.ig
CHAPTER 20
B U I L D I N G   T H E   I T E M   F I L E
.pp
Once the system parameters have been set, to build the item file, select the
ITEM ENTRY program from the system menu.  If no valid records are on the
item file, the program goes automatically into the adding mode with the
cursor waiting at the item number field.  If there are valid records on
the item file the cursor goes to the access key field; the command
CTRL-A moves the cursor to the item number field and puts the program in
the adding mode.
.he Operations: Chapter 20
.pp
To add an item to the file, give the item number and hit RETURN.  The item
number may not be more than 10 characters long, may contain only the
letters A through Z, the numbers 0 through 9, dashes, or slashes.  After
the program displays the relative record number of the item in the
access key field, enter the rest of the information for the item.
To advance the cursor
after each entry hit RETURN (if the field is not filled), to move back
one field use the command CTRL-B, to back up to make an erasure use a
CTRL-H command, and to cancel a transaction (such as an item entry),
use the command CTRL-X.
.cp15
.pp
Note the allowed entries for each field:
.sp
.li
	  DESCRIPTION:	     30 characters maximum
	  STOCKROOM QTY:     999,999
	  REORDER POINT:     999,999
	  REORDER QTY:	     999,999
	  RETAIL PRICE:      99,999.99
	  INVENTORY VALUE:   9,999,999.99
	  VENDOR NUMBER:     any 5 digit number
	  COST/UNIT:	     99,999.99
	  QTY ON ORDER:      999,999
	  DATE ORDERED:      MM/DD/YY or DD/MM/YY
	  DATE DUE:	     MM/DD/YY or DD/MM/YY
	  BACKORDER QTY:     999,999
.pp
If you wish the ACTIVITY REPORT to reflect original additions to the
file, do not make entries for the stockroom quantity, inventory value
(unless you have chosen not to use average valuation), or
cost per unit.	Instead, make all other entries for the item, then
use the ADD TO STOCK program to enter the stockroom quantity, inventory
value, and cost per unit.
.bp
.pp
When you enter an item out of order, the item file is renamed to a
filetype with a zero in the first position (e.g., IYPRT.001)
that indicates it is unsorted.	Before you can operate any programs
other than ITEM ENTRY, ITEM FILE PRINT, or ITEM FILE SORT, you must sort
the file (see Chapter 22).
.pp
Once you enter a valid item number, ESCAPE at any point adds the record
to the end of the file and readies the program for the next item.  Another
ESCAPE right after the first takes you out of the adding mode and moves
the cursor to the access key field where you can issue commands to
access records, delete records, or save records to disk (see Chapter 21).
ESCAPE from the access key field returns you to the system menu.
.pp
While entering items onto the item file, the CTRL-S command updates the
disk directory, thereby
saving the records you have created in case a catastrophic
hardware failure occurs that causes the data to be lost (such as
a pulled plug).  Whenever you have entered more records than you would
care to enter again, leave the adding mode and give the command CTRL-S.
One unpleasant experience is usually enough to demonstrate the value of this
function.
.he
.bp
.ig
CHAPTER 21
ADDING, CHANGING, DELETING, AND EXAMINING FILE ITEMS
.pp
Use the ITEM ENTRY program to add new items to the file, to change
field values for file items, to delete items no longer stocked, or
to examine items on the file.
.he Operations: Chapter 21
.pp
To add a new item select ITEM ENTRY from the system menu, and if valid
records exist on the file, issue the CTRL-A command to add records.  See
Chapter 20 for how to add an item to the file.	If adding the item to the
end of the file puts the file out of order (for example, if you
add item number A2-4600 to an item file whose last record holds item number
A3-4550) the file will have to be sorted. A disordered file is renamed to
indicate that it is out of order (from IYPRT.101 to IYPRT.001 for example).
.pp
When you add a new item to the file and
want the stockroom quantity to be reflected on the ACTIVITY REPORT or
audit proofs, do not enter that quantity through ITEM ENTRY;
instead, enter that information through the ADD TO STOCK
program.  If you've chosen inventory valuation, or want the addition
reflected on the audit proof, do not enter the stockroom quantity, unit
cost, or inventory value through ITEM ENTRY, but rather through the
ADD TO STOCK program.
.pp
To change an item already on the file, access the record by one of three
methods when the cursor is at the access key field:
.sp
.in10
.ti-4
1.  Enter the item number.
.sp
.ti-4
2.  Enter the record number relative to the beginning of the file,
in the form @Dnnn, where D is the drive reference, and nnn is the relative
record number of the item on that drive.  For example, @C25 gets the
twenty-fifth record on the item file on drive C.
.sp
.ti-4
3.  Enter the CTRL-N command to bring up the first record on the item file,
CTRL-N for the next record, and CTRL-B for the previous record.
.sp
.qi
.ad
Move the cursor forward to the desired field with RETURN, or back a
field with CTRL-B (once a record is accessed).
Type in the new value and hit RETURN to enter it.
If you want auditability, do not change the stockroom quantity, inventory
value, or cost per unit fields through the ITEM ENTRY program.	If
you choose the average valuation feature, do not change the stockroom
quantity, inventory value, or cost per unit through ITEM ENTRY.
You may alter the item number, but note that this may put the file out of
order and necessitate a sort before you can use any program besides ITEM
ENTRY, ITEM FILE PRINT, or ITEM FILE SORT.
.pp
To delete an item from the file, access the record as explained above
for changing a record and issue the command CTRL-D.  Although not physically
removed from the file, the item can no longer be accessed.  To physically
remove deleted records you must run the ITEM FILE RECOPY program (see Chapter
23). Because the record for each inventory item contains more information
than is shown on the ITEM ENTRY screen, you should delete rather than change
an item (since you cannot change all the information on the record).
.pp
To examine an inventory item, call the ITEM ENTRY program from the system
menu, and with the cursor at the access key field call the item by item number,
relative record number, or the next-record and previous-record
commands.  If you have
chosen to audit additions and depletions, it is usually best to use ITEM ENTRY
to examine the item; besides the
additional information, the audit file does not
need to  be proofed and erased, for example, if it is of the depletions type
and you want information from the additions screen.
.he
.bp
.ig
CHAPTER 22
S O R T I N G	T H E	I T E M   F I L E
.pp5
Sorting procedures vary depending on the number of item files, the
size of the item files, and the number of disk drives available.
This chapter suggests sorting procedures for each of the configurations
suggested in Chapter 19. Other procedures besides those given in this
chapter are available; you may want to experiment some to find the one
best for you.
.pp
Only item file #1 can be sorted from the system menu, and then only
if there is space available on any drive other than the input drive
to hold the workfile.  The workfiles created (and erased at the end of
the sort) equal the size of the input file.  Space to receive the sorted
output need not be available since you can instruct the system to prompt
you to remove the disk with the original input file and insert an empty
disk to receive the output file.
.pp
If space is not available for the workfile (use the CP/M STAT command;
see the CP/M Features and Facilities manual) you need to create a special
sorting disk that contains the program QSORT.COM and the necessary
sort parameter files.  When you sort the item files outside of the Inventory
Control System this special sorting disk holds the temporary workfiles.
The input file should not be so large that it cannot fit on the sorting disk
with QSORT and the small sort parameter files.
.pp
All item files have the same filename of IYPRT.  To determine whether an
IYPRT file is the first, second, or third item file, use the CP/M TYPE
command.  The first record to appear is the "header record".  The header
record contains the file's drive reference and the higest item number on
that file.
.pp
All item files have a three digit filetype (or suffix).  A zero as the
leftmost digit indicates the file is unsorted, a 1 (one) indicates it
is sorted. A 1 (one) as the middle digit indicates the file has been
sorted, but not checked for duplicates (if, of course, the first digit is
1).  Under normal circumstances, the rightmost digit appears as 1.
.mt8
.bp
.mb5
.ce
SUGGESTED SORTING PROCEDURES
.sp2
1. ONE ITEM FILE ON A 2 DRIVE SYSTEM (WITH FILE LESS THAN 1/3 DISK):
.sp6
.li
	      A 		  B

	   INT FILES	       ITEM FILE #1
	   QSORT.COM
	   PARAM FILES
	   SRT FILES
	   AUDIT FILE
.sp
.li
Sort Parameters:

	  INPUT DRIVE:	     B
	  OUTPUT DRIVE:      B
	  WORKFILE DRIVE:    B
	  SWITCH?:	     NO
.br
.sp
.in10
.ti-10
Suggested Procedure:
Run directly from menu when sorting required. The (unsorted) input file is
erased if the ITEM ENTRY program is run at any time after the sort.
.qi
.sp2
.br
2. ONE ITEM FILE ON A 2 DRIVE SYSTEM (WITH SPACE ON A FOR THE WORKFILE AND
THE FILE NO LARGER THAN 1/2 DISK):
.sp6
.li
	      A 		     B

	   INT FILES		  ITEM FILE #1
	   QSORT
	   PARAM FILES
	   SRT FILE
	   AUDIT FILES
.sp
.li
Sort Parameters:

	  INPUT DRIVE:	    B
	  OUTPUT DRIVE:     B
	  WORKFILE DRIVE:   A
	  SWITCH?:	    NO
.mt9
.bp
.mb7
.in10
.ti-10
Sorting Procedure:
Run directly from the menu when required.  The input file is erased if the
ITEM ENTRY program is run at any time after the sort.
.qi
.sp2
.br
3. ONE ITEM FILE ON A 2 DRIVE SYSTEM (WITH ITEM FILE LARGER THAN 1/2 DISK
AND INSUFFICIENT SPACE ON A FOR THE WORKFILE):
.sp6
.li
	      A 		   B

	   INT FILES		ITEM FILE #1
	   QSORT
	   PARAM FILES
	   SRT FILES
	   AUDIT FILE
.sp
.li
Sort Parameters:

	  INPUT DRIVE:	      B
	  OUTPUT DRIVE:       B
	  WORKFILE DRIVE:     A
	  SWITCH?:	      YES
.sp
.in10
.ti-10
Sorting Procedures:  Must run this sort outside the Inventory Control System.
You will need to make up a special sorting disk with QSORT.COM, the
sort parameter file PART1.SRT, and the CP/M operating system.  Perform the
following steps:
.sp
1. To make the special sorting disk, end the Inventory Control System and
remove the disk currently on drive B.  Insert an empty, formatted disk
(with the CP/M operating system) on
B.  Give the CTRL-C to reset the maps of disk usage.   Type the command
PIP B:=A:QSORT.COM[V], the command PIP B:=A:*.SRT[V], and the command
PIP B:=A:PIP.COM[V] (for convenience).	The B disk is now your special
sorting disk.
.bp
2. Remove the disk on drive A and replace with the sorting disk on B.
Put the disk that was originally on B on drive B. Give the CTRL-C command.
The A drive should now hold
the following files: QSORT.COM, PART1.SRT, and PIP.COM.  The B drive
should hold only the first item file, IYPRT.001 (the filetype 001 indicates
the item file is unsorted; when sorted it changes to type 101).
.sp
3. Type the command QSORT PART1.  When the sort finishes, the system prompts
you to remove the B disk with the unsorted item file, and replace it with
an empty disk to receive the sorted output from the workfiles on A.  Do not
reboot.
.qi
.sp2
.br
4. TWO ITEM FILES ON A 2 DRIVE SYSTEM:
.sp6
.li
	      A 		 B

	   INT FILES	      ITEM FILE #2
	   QSORT	      PARAM FILES
	   ITEM FILE #1       SRT FILES
			      AUDIT FILE
.sp
.li
Sort Parameters for Item File #1:

	  INPUT DRIVE:	    A
	  OUTPUT DRIVE:     A
	  WORKFILE DRIVE:   B
	  SWITCH?:	    YES

Sort Parameters for Item File #2:

	  INPUT DRIVE:	    B
	  OUTPUT DRIVE:     B
	  WORKFILE DRIVE:   A
	  SWITCH?:	    YES
.sp
.in10
.ti-10
Sorting Procedure:  Item file #1 can be sorted from the menu, without
disk switching if there is room on A to receive the sorted output, and room
on B to hold the temporary workfiles (which are the same size as the input
file).	Otherwise, sort both files outside the Inventory Control System
using the parameters given above.
.sp
1. Prepare a special sorting disk containing QSORT, PIP, and the two SRT files
(and the CP/M operating system).   In this case you need to get QSORT from
the disk that was on A, and the SRT files from the disk that was on B
(the "system drive").
.sp
2. Put the sorting disk on B, and the disk with item file #1 on A.  Give
the CTRL-C command.  Type the command B:QSORT PART1 (to run the sort from
drive B).  Switch the A disk to receive the sorted output when prompted.
.sp
3. Remove the disk now in A, set it aside, and insert the special sorting
disk that was on B into A.  Put the disk with item file #2 on drive B.
Give the CTRL-C command.  Type the command QSORT PART2.  When instructed
by the program remove the B disk and replace it with a disk to
receive the sorted output from A.
.sp
4. You now have the two original disks which contain the unsorted item files
(and other files), one sorting disk, and two disks containing only the sorted
files.	You must now either PIP the files that were on the original disks
onto the new, sorted, disks, or erase the unsorted files from the old
disks and PIP the sorted files back in their place.  Make sure the files
are allocated properly before re-entering the Inventory Control System.
.qi
.sp2
.br
5. ONE ITEM FILE ON A 3 DRIVE SYSTEM:
.sp6
.li
	      A 		   B			C

	   INT FILES		PARAM FILES	     ITEM FILE #1
	   QSORT		SRT FILES
				AUDIT FILE

.bp
.li
Sort Parameters for Item File #1:

	  INPUT DRIVE:	     C
	  OUTPUT DRIVE:      C
	  WORKFILE DRIVE:    B
	  SWITCH?:     YES (if file is greater than 1/2 disk)
		       NO  (if file is less than 1/2 disk)
.sp
.in10
.ti-10
Sorting Procedure:  May be run from the system menu.  Switch disks when
prompted if switching requested.  If the parameter files, and audit file on
drive B occupy too much space and do not allow the workfile to fit, you
should PIP those files onto the A drive and erase them from the B drive.
Or, you could prepare the special sorting disk as explained above and
run the sort outside the system.
.qi
.sp2
.br
6. TWO ITEM FILES ON A 3 DRIVE SYSTEM (WITH SPACE AVAILABLE ON A AT LEAST
EQUAL TO THE SIZE OF ITEM FILE #1):
.sp6
.li
	      A 	       B	      C

	   INT FILES	    ITEM FILE #1   ITEM FILE #2
	   QSORT
	   PARAM FILES
	   SRT PARAMS
	   AUDIT FILE

Sort Parameters for Item File #1:

	  INPUT DRIVE:	     B
	  OUTPUT DRIVE:      B
	  WORKFILE DRIVE:    A
	  SWITCH?:     YES (if file is larger than 1/2 disk)
		       NO  (if file is less than 1/2 disk)

Sort Parameters for Item File #2:

	 INPUT DRIVE:	    C
	 OUTPUT DRIVE:	    C
	 WORKFILE DRIVE:    B
	 SWITCH?:     YES (if file is greater than 1/2 disk)
		      NO  (if file is less than 1/2 disk)
.br
.bp
.in10
.ti-10
Sorting Procedure:  May sort item file #1
directly from the system menu.	Be sure
QSORT is on the currently logged drive (which in this case should be drive
A), and PART1.SRT is on the drive named as the system drive (also drive
A for this configuration).
.sp
1. Select the SORT ITEM FILE program from the menu and switch disks when
prompted.
.sp
2. Sort Item File #2  outside of the system.  Follow the procedure for sorting
item file #2 as explained in the section below.
.qi
.sp2
.br
7.  TWO ITEM FILES ON A 3 DRIVE SYSTEM (WITH INSUFFICIENT SPACE AVAILABLE
ON DRIVE A FOR THE ITEM FILE #1 WORKFILE):
.sp6
.li
	      A 		 B		    C

	   INT FILES	      ITEM FILE #1	 ITEM FILE #2
	   QSORT
	   PARAM FILES
	   SRT FILES
	   AUDIT FILE

Sort Parameters for Item File #1:

	  INPUT DRIVE:	     B
	  OUTPUT DRIVE:      B
	  WORKFILE DRIVE:    C
	  SWITCH?:     YES (if file is larger than 1/2 disk)
		       NO  (if file is less than 1/2 disk)

Sort Parameter for Item File #2:

	  INPUT DRIVE:	     C
	  OUTPUT DRIVE:      C
	  WORKFILE DRIVE:    B
	  SWITCH?:     YES (if file is larger than 1/2 disk)
		       NO  (if file is less than 1/2 disk)
.bp
.in10
.ti-10
Sorting Procedure:  Must sort both item files outside the system.
Be sure QSORT is on the currently logged drive, and the sort parameter
files on the drive named as the system drive (these should both be drive
A under this configuration).
.sp
1. With QSORT, PART1.SRT, and PART2.SRT on
drive A, put the disk with item file #1 on drive B.  Put an empty, formatted,
disk on drive C.  Give the CTRL-C command.
.sp
2. Type QSORT PART1.  When the sort routine finishes, switch the B disks
to receive the sorted output if switching was ordered.
.sp
3. Remove the disk now in B and replace with the blank disk that served as
the workfile on C.  Put the disk with item file #2 on drive C and give
the CTRL-C command.
.sp
4. Type QSORT PART2.  Switch disks on drive C to receive the sorted output
if switching was ordered.  Remove the blank disk from drive B, replace
it with the sorted item file #1, give the CTRL-C command, and you're ready
to re-enter the Inventory Control System.
.qi
.sp2
.br
8.  THREE ITEM FILES ON A 3 DRIVE SYSTEM:
.sp6
.li
	      A 		  B		     C

	   INT FILES	       ITEM FILE #2	  ITEM FILE #3
	   QSORT	       PARAM FILES
	   ITEM FILE #1        SRT FILES
			       AUDIT FILE

.bp
.li
Sort Parameters for Item File #1:

	  INPUT DRIVE:	     A
	  OUTPUT DRIVE:      A
	  WORKFILE DRIVE:    B
	  SWITCH?: YES (if space on A less than item file #1)
		   NO  (if space on A greater than item file #1)

Sort Parameters for Item File #2:

	  INPUT DRIVE:	     B
	  OUTPUT DRIVE:      B
	  WORKFILE DRIVE:    A
	  SWITCH?: YES (if space on B less than item file #2)
		   NO  (if space on B greater than item file #2)

Sort Parameters for Item File #3:

	  INPUT DRIVE:	     C
	  OUTPUT DRIVE:      C
	  WORKFILE DRIVE:    A
	  SWITCH?: YES (if space on C less than item file #3)
		   NO  (if space on C greater than item file #3)
.br
.sp
Sorting Procedure: All files must be sorted outside the system.
.sp
.in10
1.  Make a special sorting disk that contains the CP/M system, QSORT, PIP,
and the three
sort parameter files PART1.SRT, PART2.SRT, and PART3.SRT
Place this
disk on drive B.  Put the disk with item file #1 on drive A and
give the CTRL-C command.
.sp
2.  Type B:QSORT PART1.  Switch A disks to receive the sorted output
from the workfile on B if switching has been requested.
.sp
3.  To sort item file #2, put the disk that contains that file on drive
B, the sorting disk on A, and give the CTRL-C command.	Type
QSORT PART2.  Switch B disks to receive the sorted output from the workfile
on A if switching has been requested.
.sp
4.  To sort item file #3, put the disk that contains that item file on drive
C, the sorting disk on drive A, and give the CTRL-C command.  Type
QSORT PART3.  Switch C disks to receive the sorted output from the workfile
on A if switching has been requested.
.sp
5.  You now have the original A, B, and C disk which contain the unsorted
item files and the other Inventory Control System programs and files, one
special sorting disk, and three disks containing only the sorted item files.
To get back to the original configuration, you must either PIP the programs
and other files onto the disks with the sorted files, or
erase the unsorted files from the original
disks, and PIP the sorted files onto the original disks in their place.
Be sure all programs and files reside on the proper disks before restarting
the Inventory Control System.
.qi
.sp2
.br
9.  ONE ITEM FILE ON A 4 DRIVE SYSTEM
.sp6
.li
	   A		  B		    C		   D

	INT FILES      PARAM FILES	 ITEM FILE #1	not used, or
	QSORT	       SRT FILES			AUDIT FILE
		       AUDIT FILE

Sort Parameter for Item File #1:

	  INPUT DRIVE:	     C
	  OUTPUT DRIVE:      C
	  WORKFILE DRIVE:    A, B, or D
	  SWITCH?:     YES (if file is larger than 1/2 disk)
		       NO  (if file is less than 1/2 disk)
.br
.sp
.in10
.ti-10
Sorting Procedure:  Run directly from the system menu.	Be sure QSORT is
on the currently logged drive and the SRT files on the drive named as
the system drive (in this case, B).
.qi
.bp
10.  TWO ITEM FILES ON A 4 DRIVE SYSTEM
.sp6
.li
      A 	       B		   C		    D

   INT FILES	    PARAM FILES 	ITEM FILE #1	 ITEM FILE #2
   QSORT	    SRT FILES
		    AUDIT FILE

Sort Parameters for Item File #1:

	  INPUT DRIVE:	     C
	  OUTPUT DRIVE:      C
	  WORKFILE DRIVE:    D
	  SWITCH?:     YES (if file is larger than 1/2 disk)
		       NO  (if file is less than 1/2 disk)

Sort Parameters for Item File #2:

	  INPUT DRIVE:	     D
	  OUTPUT DRIVE:      D
	  WORKFILE DRIVE:    C
	  SWITCH?:     YES (if file is larger than 1/2 disk)
		       NO  (if file is less than 1/2 disk)
.br
.sp
Sorting Procedure:
.sp
.in10
1. To sort item file #1,
be sure that QSORT is on the currently logged drive (in this case
drive A) and the sort parameter files on the system drive (in this case
drive B).  Put the disk with item file #1 on drive C, and a empty, formatted,
disk on drive D.  Give the CTRL-C command.
Type QSORT PART1.  Switch the C disks when prompted if
switching has been requested.
.bp
2. To sort item file #2,
make sure QSORT is on the currently logged drive (A) and the
sort parameter (SRT) files on the
system drive (B). Put an empty, formatted, disk on drive B, and the unsorted
item file #2 on drive D.  Give the CTRL-C command.  Type QSORT PART2.  Switch
D disks to receive the sorted output from the workfile if switching has been
requested.
.sp
Make sure all disks and files are in their correct locations before beginning.
.qi
.sp2
11.  THREE ITEM FILES ON A 4 DRIVE SYSTEM:
.sp6
.li
	   A		     B		   C		    D

	INT FILES	  ITEM FILE #1	ITEM FILE #2	 ITEM FILE #3
	QSORT
	PARAM FILES
	SRT FILES
	AUDIT FILE

Sort Parameters for Item File #1:

	  INPUT DRIVE:	     B
	  OUTPUT DRIVE:      B
	  WORKFILE DRIVE:    A
	  SWITCH?:     YES (if file larger than 1/2 disk)
		       NO  (if file is less than 1/2 disk)

Sort Parameters for Item File #2:

	  INPUT DRIVE:	     C
	  OUTPUT DRIVE:      C
	  WORKFILE DRIVE:    A
	  SWITCH?:     YES (if file larger than 1/2 disk)
		       NO  (if file is less than 1/2 disk)

.bp
.li
Sort Parameters for Item File #3:

	  INPUT DRIVE:	     D
	  OUTPUT DRIVE:      D
	  WORKFILE DRIVE:    A
	  SWITCH?:     YES (if file larger than 1/2 disk)
		       NO  (if file is less than 1/2 disk)
.br
.sp
Sorting Procedure:
.sp
.in10
1.  Make up a special sorting disk containing the CP/M system, QSORT, PIP,
and the three sort parameter
files (SRT).
.sp
2.  Put the special sorting disk in drive A, the disk
with item file #1 on drive B, the disk with item file #2 on drive C, and
the disk with item file #3 on drive D.	Give the CTRL-C command.
.sp
3.  Type QSORT PART1. Switch B disks when prompted if switching requested.
.sp
4.  Type QSORT PART2. Switch C disks when prompted if switching requested.
.sp
5.  Type QSORT PART3.  Switch D disks when prompted if switching requested.
.sp
6.  Replace the sorting disk in A with the disk containing the other programs
and files.  Give the CTRL-C command before restarting the Inventory Control
System.
.qi
.he
.bp
.ig
CHAPTER 23
R E C O P Y I N G    T H E    I T E M	F I L E
.pp5
When you alter the item file parameters (see Chapter 19), or when you want
to remove deleted records from the file, select the ITEM FILE
RECOPY program from the system menu.  The program first asks if you want
to change the item file parameters.  YES allows you to do so.  NO
tells the program you just want to remove deleted records.
.he Operations: Chapter 23
.pp
Under the direction and prompting of the ITEM FILE RECOPY program, each
item file is recopied, in turn, from the A drive to the B drive.  When
you are instructed, put the disk with the item file to be recopied
on drive A and
a new disk on B.  Hit RETURN to begin the recopy.
.pp
The program tells you when to remove the disk from the A drive
and replace it with the disk containing the
next item file to be recopied.	The program
also prompts you to replace the disk on the B drive with another disk
at the proper time. The disk you place on B should never contain another item
file (a file with the filename of IYPRT).
.pp
Since the disks with the item files being recopied often contain other
files as well, these files must be moved to the new disks
at some point before you can operate the Inventory Control System.
This is best
done before starting the recopy program, by using the PIP command.
Be sure there is room enough on each disk to receive the recopied files.
For example, you would not want to increase the size of an item file
(by changing the item file parameters) so
much that it does not fit on the disk with the other files.  If this
happens, you should re-allocate the Inventory Control System programs and
files (see Chapter 19 for suggestions).
.pp
After all the item files have been recopied (you must recopy all of them when
running the recopy program), the program checks to see if your system disk
is in the drive you designated as the system drive.  The system disk is
the disk with the two parameter files (IYPR1 AND IYPR2) and the
sort parameter files, not necessarily the disk with the INT files.  If
it does not find the system disk where it should be, the program instructs
you to put it there.  After checking for the system disk, the program
looks on the currently logged drive for the disk with the INT files, and
prompts you to put it there if it is not already.
.he
.bp
.ig
CHAPTER 24
A D D I N G   T O   S T O C K
.pp5
Use the ADD TO STOCK program to increase the stockroom quantity, or
change order information for an inventory item.
.he Operations: Chapter 24
.pp
To use this program, the item files must be sorted by item number
(see Chapter 22), and the audit file must be of the additions type
if you have requested additions transactions to be audited.  If the
audit file is not an additions audit file, you should print a transaction
proof of the current depletions file, and erase it,
before you begin the add to stock
transactions.  The ADD TO STOCK program creates a new additions audit file
and increases the "batch number" by one.
.pp
When the program asks if you want to erase the audit file, you may respond
with Y (YES), which erases it, with N (NO), which leaves it intact and
ends the ADD TO STOCK program, or with ESCAPE, which is the same as NO.
.pp
To make an addition to stock, first call up the record for that inventory
item.  With the cursor at the ITEM NUMBER field,
.sp
.in10
.ti-4
1.  Enter the item number in its exact form.
.sp
.ti-4
2.  Enter an @ (at-sign) followed by the drive reference and record number
relative to the beginning of the item file on that drive.  For example,
to get the 36th record on item file #2 located
on drive B, enter "@B36" followed
by RETURN.  The @ (at-sign) indicates that what follows it is a relative
record number, and not an item number.
.sp
.ti-4
3.  Give the CTRL-N command to bring up the first record on the first
item file, and after that, to bring up the record following whichever
record is present. CTRL-B brings up the record immediately preceding
the record currently displayed.
.qi
.sp
.ad
Enter the quantity added to stock and hit RETURN to advance to the next
field if the quantity added field is not filled.  Enter the cost per unit
of the stock received and update the order information fields.	Do not enter
commas, the program inserts them where needed. If an order
is filled and you wish to clear the date fields, type "none" in the field.
You may cycle through the fields with RETURN and CTRL-B (for going back
a field), adjusting the information as necessary until it is correct.
The record itself, and the associated accumulators are not updated until
you hit ESCAPE.  You may then enter another addition to stock transaction,
or enter ESCAPE again to return to the system menu.
.pp
An addition to stock increases the stockroom quantity and decreases the
quantity on order and the quantity back-ordered.  The quantity back-ordered
is treated as a subset of the quantity on order.
.pp
To back out (correct) an addition to stock that was made too large,
enter a negative
number in the amount of the error (at the same unit cost as the original
entry).  Do not correct such an error by an offsetting depletion from stock as
this will make the activity report inaccurate.	For example, if you discover
you have made an entry too large by 10 units, enter -10 in the quantity
added field and the appropriate value in the cost per unit field
to correct the error.  If the entry was made too small, just add the
additional units in the regualar way.
.he
.bp
.ig
CHAPTER 25
D E P L E T I N G    S T O C K
.pp5
Use the DEPLETE STOCK program to decrease the stockroom quantity
when goods leave your inventory.
.he Operations: Chapter 25
.pp
To use this program, the item files must be sorted by item number
(see Chapter 22), and the audit file must be of the depletions type
if you have requested depletions transactions to be audited.  If
the audit file is not a depletions audit file, you should print
a transaction proof of the current additions file, and erase it, before
you begin the deplete stock transactions.  The DEPLETE STOCK program
creates a new depletions audit file and increases the "batch number"
by one.
.pp
When the program asks if you want to erase the sudit file you may respond
with Y (YES), which erases it, with N (NO), which ends the DEPLETE STOCK
program, or with ESCAPE, which is the same as NO.
.pp
To make a depletion from stock, first call up the record for that inventory
item.  With the cursor at the item number field,
.sp
.in10
.ti-4
1.  Enter the item number in the exact form.
.sp
.ti-4
2.  Enter an @ (at-sign) followed by the drive reference and record number
relative to the beginning of the item file on that drive.  For example,
to get the 36th record on item file #2 located on drive B, enter
"@B36" followed by RETURN.  The @ (at-sign) indicates that what follows it
is a relative record number, and not an item number.
.sp
.ti-4
3.  Give the CTRL-N command to bring up the first record on the first item
file, and after that, to bring up the record following whichever record
is present.  CTRL-B brings up the record immediately preceding the record
currently displayed.
.qi
.sp
.ad
Enter the quantity leaving stock and follow with RETURN.
The program updates the stockroom
quantity, and, if average valuation has been requested, decreases the
inventory value based on the current average value the item.  The cursor
returns to the item number field where you can enter a new number to begin
the next transaction.
ESCAPE returns you to the system menu.
.pp
To back out (correct) a depletion made too large, enter a negative
amount in the quantity depleted field in the amount of the error.  Do not
correct the error by making an addition to stock, as this makes the
activity report inaccurate.
.bp
.pp
If you've chosen the average valuation feature, an addition to stock
between the time a depletion error is made and the time it is discovered
may change the average value of that item.  Consequently, the adjustement
of the error corrects the stockroom quantity, but the new inventory
value, computed using a different average value for the item, may not be the
same.  The difference will usually not be critical, however, since
inventory value is an estimate based on averages anyway.  Allowing only
an addition or depletion audit file to exist at one time helps prevent this
situation from arising, since a careful examination of the transaction proofs
may allow you to
catch the error before an addition to stock changes the average value
of the inventory item.
.he
.bp
.ig
CHAPTER 26
P R O D U C I N G   T H E   R E P O R T S
.pp5
Once system setup is complete you can print Inventory Control
System reports from the system menu at any time.  For all but the
ITEM FILE PRINT, the item files must be in sorted order.
.pp
To print a report, select the desired report program from the system menu
and choose the report options you wish to employ.
Be sure your printer is ready and the paper aligned.
After the printout finishes, the system menu returns.
.pp
For certain reports, a range option is available that lets you select
a range of high and low item numbers.  Leaving the range
values blank by hitting RETURN reports on all inventory items. Giving
a low value but not a high value reports on all items with numbers
above and including that value.  Giving a high value but not a low value
reports on all items with numbers below and including that value.
.pp
If the average valuation option has not been chosen, the VALUATION
REPORT will not be useful unless you have maintained the inventory
value manually.
.pp
At the end of each current period or to-date period, reset the
appropriate activity registers to zero by running the START AN ACTIVITY
PERIOD program and answering its requests.
#
